\bookchapter{Глава 17: Связь телефона с другими устройствами}

\chapterintro{Ваш телефон — это на самом деле супер-пульт и центр данных. Он не только звонит, но и может объединить компьютер, телевизор, часы, колонки и автомобиль в единую сеть, чтобы эти устройства перестали быть изолированными островами. В этой главе мы разберём, как заставить их «работать в команде».}

\booksection{I. Связка с компьютером: телефон и компьютер больше не два разных мира}

Честно говоря, раньше меня больше всего раздражало, когда я снимал фото на телефон, хотел обработать его на компьютере, а приходилось искать кабель и ждать передачи. Потом я понял, что кабель вообще не нужен — хорошая связка работает в десять раз быстрее.

\begin{itemize}
  \item \textbf{Экосистема Apple (AirDrop + Универсальный буфер обмена)}: если у вас iPhone и Mac, AirDrop — первая функция, которую стоит освоить. На телефоне нажмите «Поделиться» AirDrop выберите ваш Mac, и фото мгновенно перенесётся. Ни интернет, ни Wi-Fi не нужны — используется Bluetooth и прямое Wi-Fi-соединение. Ещё одна легко упускаемая функция — «Универсальный буфер обмена»: скопируйте текст на телефоне и вставьте на Mac — работает как магия.
  \item \textbf{Android + Windows (Связь с Windows)}: у Microsoft есть приложение «Связь с Windows», встроенное в телефоны Samsung; для других Android-телефонов достаточно скачать его из магазина приложений. После сопряжения вы можете на компьютере просматривать уведомления телефона, отвечать на SMS, листать фото и даже принимать звонки. Я пользуюсь этим два года — на совещаниях не нужно постоянно опускать глаза к телефону, достаточно уведомления на экране компьютера.
  \item \textbf{Фирменная мультиэкранная связка}: у Huawei «Мультиэкранная связка» и у Xiaomi «Xiaomi Share» ещё круче — экран телефона напрямую транслируется на компьютер, можно перетаскивать файлы мышью с телефона на компьютер, а также смотреть TikTok и отвечать в WeChat на компьютере. Кроссплатформенное копирование и вставка, перетаскивание файлов — всё работает очень плавно.
\end{itemize}

\booksection{II. Трансляция на телевизор: смотрим всей семьёй}

Когда я приезжал домой на Новый год, моя кузина хотела показать на телевизоре видео с её выпускной поездки. Долго мучились, а потом оказалось, что всё очень просто:

\begin{itemize}
  \item \textbf{iPhone Телевизор}: смахните вниз из правого верхнего угла для Центра управления, нажмите значок двух наложенных прямоугольников (Повтор экрана), найдите название вашего телевизора — и всё подключено. Условие: телевизор и телефон в одной сети Wi-Fi, и телевизор поддерживает AirPlay (большинство современных Smart TV поддерживают).
  \item \textbf{Android Телевизор}: у каждого бренда своё название: Huawei — «Беспроводная трансляция», Xiaomi — «Трансляция», OPPO/vivo — «Мультиэкранное взаимодействие». Обычно находится в панели быстрых настроек; если нет — поищите «Трансляция» в настройках.
  \item \textbf{Прямая трансляция из видеоприложений}: если вы просто хотите посмотреть iQiyi, Bilibili, YouTube на телевизоре, всё ещё проще — нажмите на значок прямоугольника с символом Wi-Fi в интерфейсе воспроизведения, выберите телевизор, и видео пойдёт на него. Телефон в это время можно использовать для других дел — видео продолжает играть на телевизоре.
  \item \textbf{Что делать, если трансляция не работает}: в 90\% случаев проблема в том, что устройства не в одной сети Wi-Fi. В оставшихся 10\% — функция трансляции не включена на телевизоре. Поищите в настройках телевизора «Мультиэкранное взаимодействие» или «AirPlay».
\end{itemize}

\booksection{III. Умные часы и браслеты: монитор здоровья на запястье}

Я купил маме фитнес-браслет. Сначала она сказала: «Зачем мне, я же не занимаюсь спортом». А через две недели обнаружила, что спит всего чуть больше пяти часов в сутки, и начала сама корректировать режим. В этом самая практичная ценность часов и браслетов — превращать невидимые данные о здоровье в видимые цифры.

\begin{itemize}
  \item \textbf{Фильтр уведомлений}: когда телефон вибрирует без остановки в сумке или кармане, поднимите запястье — и за секунду поймёте, это начальник в WeChat или реклама Pinduoduo. Сообщения, не стоящие того, чтобы доставать телефон, пусть идут мимо.
  \item \textbf{Синхронизация данных о здоровье}: часы круглосуточно измеряют пульс, кислород в крови, сон, шаги — эти данные автоматически собираются в приложении «Здоровье» на телефоне. Со временем вы увидите тренды: достаточно ли вы ходите в этом месяце, улучшилось ли качество сна на этой неделе.
  \item \textbf{Синхронизация транспортных карт и пропусков}: привязанные на телефоне транспортные карты и ключи-пропуска большинство брендов автоматически синхронизируют с часами. В метро не нужно доставать телефон — приложили запястье к турникету и прошли. Реально возникает ощущение «будущее уже здесь».
  \item \textbf{Поиск телефона}: на часах обычно есть функция «Найти телефон» — нажали, и телефон издаёт звук. Я пользуюсь этой функцией гораздо чаще, чем ожидал — вечно кладу телефон куда-то и забываю.
\end{itemize}

\booksection{IV. Умный дом: телефон как «главный выключатель»}

Умный дом звучит загадочно, но на самом деле сводится к трём вещам: удалённое управление с телефона, настройка автоматических сценариев и голосовое управление.

Расскажу на примерах из собственной жизни, и вы поймёте, насколько это практично:

\begin{itemize}
  \item \textbf{Кондиционер}: летом по дороге домой заранее включаю кондиционер с телефона — захожу в уже прохладную квартиру. За сто с лишним юаней можно купить умную розетку-переходник для кондиционера, менять сам кондиционер не нужно.
  \item \textbf{Освещение}: лёжа в кровати перед сном, лень вставать выключать свет — нажал в телефоне или сказал: «Привет, Siri / Xiao Ai, выключи свет».
  \item \textbf{Робот-пылесос}: перед уходом нажал в телефоне «Начать уборку» — вернулся, а пол чистый. Можно настроить ежедневную уборку по расписанию.
  \item \textbf{Камеры}: если дома пожилые люди, дети или питомцы — можно в любой момент посмотреть, что происходит дома, через телефон.
\end{itemize}

Какую экосистему выбрать? Краткий ответ:

\begin{itemize}
  \item iPhone выбирайте устройства с поддержкой Apple HomeKit (дороже, но с хорошей защитой приватности).
  \item Максимальное соотношение цена/качество экосистема Xiaomi Mijia самая полная: лампы, розетки, пылесосы — всё есть.
  \item Телефон Huawei выбирайте HarmonyOS Connect — лучший опыт с устройствами Huawei.
\end{itemize}

\booksection{V. Bluetooth-наушники и колонки: свобода без проводов}

Переход с проводных наушников на Bluetooth-наушники стал для меня самым заметным улучшением в опыте использования телефона. Больше никогда не достаёшь из сумки клубок проводов и не боишься зацепиться проводом за дверную ручку.

\begin{itemize}
  \item \textbf{Сопряжение очень простое}: при первом включении наушники обычно автоматически входят в режим сопряжения (мигающий индикатор). Включите Bluetooth на телефоне, найдите наушники и нажмите для подключения. В дальнейшем они подключаются автоматически при извлечении из кейса — ничего делать не нужно.
  \item \textbf{Бонусы одного бренда}: AirPods с iPhone автоматически показывают всплывающее окно с уровнем заряда и могут бесшовно переключаться между устройствами Apple. Huawei FreeBuds с телефонами Huawei, наушники Xiaomi с телефонами Xiaomi — аналогичные всплывающие окна и оптимизация.
  \item \textbf{Bluetooth-колонки}: на семейных праздниках или пикниках одна Bluetooth-колонка удваивает атмосферу. Подключили — и телефон становится вашим диджейским пультом.
  \item \textbf{Внимание к задержке}: при игре в игры с Bluetooth-наушниками может быть небольшая задержка (выстрел слышен через полсекунды). Если задержка критична, используйте проводные наушники или купите Bluetooth-наушники с поддержкой «режима низкой задержки».
\end{itemize}

\begin{tcbinfo}{UWB (Ultra-Wideband) — более точное «пространственное чутьё», чем Bluetooth}
Помимо Bluetooth, в новых телефонах скрыт чип UWB (Ultra-Wideband, сверхширокополосная связь). Его способность —\textbf{точное измерение расстояния и пространственное позиционирование}— он может определить, насколько далеко и в каком направлении находится другое устройство, с точностью до сантиметров. Bluetooth может сказать только «рядом», а UWB — «в 1,2 метра слева впереди». Распространённые применения:

\begin{itemize}
  \item \textbf{AirTag / точный поиск}: когда iPhone ищет AirTag, стрелка-навигатор показывает «7 метров, идите налево» — это работает благодаря UWB. Обычный Bluetooth может только сказать, в зоне ли действия.
  \item \textbf{Цифровой ключ автомобиля}: iPhone и флагманы Android с UWB могут работать как ключ: подходите к машине — автоматически разблокируется, отходите — запирается. Может различать положение водителя и пассажира, защита от ретрансляционных атак.
  \item \textbf{Навигация внутри помещений}: в торговых центрах и аэропортах с UWB-маяками возможна навигация с точностью до прохода — гораздо точнее GPS.
  \item \textbf{Поддерживаемые модели}: iPhone 11 и новее оснащены чипом UWB (U1/U2); среди Android — флагманы Xiaomi, Samsung, vivo, OPPO.
\end{itemize}

\end{tcbinfo}

\begin{tcbinfo}{LE Audio и Auracast — новое поколение стандартов Bluetooth-аудио}
Классический Bluetooth-аудио (A2DP), прослуживший двадцать лет, сменяется\textbf{LE Audio}(Low Energy Audio) — новым стандартом аудио начиная с Bluetooth 5.2:

\begin{itemize}
  \item \textbf{Экономичнее, качество лучше}: новый кодек LC3 обеспечивает лучшее качество звука, чем старый SBC, при меньшем энергопотреблении — автономность наушников заметно увеличивается.
  \item \textbf{Двухканальное аудио}: левый и правый наушники передают сигнал независимо, переключение в режим одного наушника без разрыва.
  \item \textbf{Auracast Broadcast Audio}: фирменная функция LE Audio. Один телефон или телевизор может одновременно транслировать аудио на\textbf{неограниченное количество}наушников — в залах ожидания аэропортов, музеях, спортзалах, классах можно использовать один источник для всех желающих.
  \item \textbf{Текущая ситуация}: новые флагманские наушники и телефоны 2024–2026 годов повсеместно поддерживают LE Audio. Старые устройства требуют аппаратного обновления. При покупке новых наушников обращайте внимание на маркировку LE Audio / Auracast.
\end{itemize}

\end{tcbinfo}

\booksection{VI. Автомобильные системы: безопасный партнёр за рулём}

Смотреть в телефон за рулём — незаконно и смертельно опасно. Но иногда действительно нужна навигация, приём звонков, музыка. Здесь на помощь приходят CarPlay (Apple) и Android Auto (Android).

\begin{itemize}
  \item \textbf{Как подключить}: если ваш автомобиль поддерживает (большинство новых машин последних лет), подключите телефон кабелем к USB-порту автомобиля — экран автоматически переключится на интерфейс CarPlay или Android Auto. Некоторые новые авто поддерживают беспроводное подключение.
  \item \textbf{Что можно делать}: навигация (гораздо лучше встроенной), звонки (приём с руля), музыка (приложения с телефона напрямую на авто), голосовые сообщения (Siri или Xiao Ai читают и отвечают).
  \item \textbf{Для старых автомобилей}: если машина старая и не поддерживает, можно купить внешний экран CarPlay за несколько сотен юаней и прикрепить к приборной панели — опыт почти такой же.
\end{itemize}

\booksection{VII. Безопасность и конфиденциальность при связке устройств}

Чем больше устройств связано, тем больше поверхность риска утечки приватности. На что обратить внимание:

\begin{itemize}
  \item Перед трансляцией экрана в публичном месте проверьте, нет ли в уведомлениях личных сообщений — не показывайте всем то, что не предназначено для чужих глаз.
  \item Перед тем как подарить или продать устройство, отвяжите его — зайдите в настройки телефона, найдите устройство и выберите «Удалить» или «Отменить сопряжение».
  \item Для камер умного дома покупайте брендовые, с шифрованием. Не берите дешёвые ноунеймы за несколько десятков юаней. Слишком много новостей о взломанных камерах с прямой трансляцией для хакеров.
\end{itemize}

\begin{tcbpractice}{Попробуйте: выведите экран телефона на телевизор и создайте умный сценарий}
\begin{enumerate}
  \item Убедитесь, что телефон и телевизор подключены к одной сети Wi-Fi.
  \item На телевизоре найдите функцию «Повтор экрана» или «AirPlay» и включите её.
  \item На телефоне: iPhone — из Центра управления нажмите «Повтор экрана»; Android — потяните вниз для «Беспроводной трансляции».
  \item Выберите любимую фотографию и выведите её на телевизор, покажите семье.
  \item Если есть устройства умного дома, попробуйте создать в приложении сценарий «Ухожу из дома»: одно нажатие — выключить свет, выключить кондиционер, запустить робот-пылесос.
\end{enumerate}

\end{tcbpractice}

\begin{tcbinfo}{Кейс 1: Практическая оптимизация автономности батареи}
\begin{itemize}
  \item \textbf{Сценарий}: пользователь Xiaomi 14 — плохая автономность, разрядка за полдня.
  \item \textbf{Шаги}: (1) Настройки Батарея посмотреть рейтинг расхода (2) Определить приложения с аномальным фоновым расходом (3) Отключить 5G, переключиться на 4G (4) Ограничить фоновую активность (5) Включить режим энергосбережения (6) Снять чехол при зарядке для теплоотвода.
  \item \textbf{Ожидаемый результат}: автономность увеличилась на 30–50\%, данные о здоровье батареи читаемы.
\end{itemize}

\end{tcbinfo}
